% =============================================================================
% COMPONENTES EDITORIAIS REUTILIZÁVEIS (PADRÃO PROFISSIONAL E DISCRETO)
% =============================================================================

\usepackage[most]{tcolorbox}
\usepackage{enumitem}
\usepackage{listings}
\usepackage{tabularx}
\usepackage{booktabs}
\usepackage{colortbl}
\usepackage{array}
\usepackage{float}
\usepackage[export]{adjustbox}

% Padronização tipográfica e alinhamento vertical/horizontal de tabelas
\renewcommand{\tabularxcolumn}[1]{m{#1}}
\newcolumntype{C}[1]{>{\centering\arraybackslash}m{#1}}
\newcolumntype{L}[1]{>{\raggedright\arraybackslash}m{#1}}
\newcolumntype{J}[1]{>{\sloppy\arraybackslash}m{#1}}
\newcolumntype{Y}{>{\sloppy\arraybackslash}X}

% Marcador de lista limpo para as caixas temáticas
\newcommand{\itemcaixa}[1]{%
  \par\vspace{5pt}\noindent
  \hangindent=0.55cm\hangafter=1
  \makebox[0.55cm][l]{\textcolor{white!80}{\fontsize{9pt}{9pt}\selectfont •}}{\color{white}#1}\par
}

\newcommand{\itemsaibamais}[1]{\itemcaixa{#1}}

\newcommand{\linkcaixa}[2]{%
  \href{#1}{\textcolor{white}{\textbf{\underline{#2}}}}%
}

% -----------------------------------------------------------------------------
% COMPONENTE 2: COMANDOS SEMÂNTICOS DE SEÇÃO
% -----------------------------------------------------------------------------
\newcommand{\secao}[1]{\section{#1}}
\newcommand{\subsecao}[1]{\subsection{#1}}
\newcommand{\subsubsecao}[1]{\subsubsection{#1}}

\newcommand{\topico}[1]{%
  \par\vspace{6pt}\noindent
  {\fontsize{10pt}{13.5pt}\fontExtraBold\color{BrandNavy}#1}\par\nopagebreak\vspace{2pt}%
}

% -----------------------------------------------------------------------------
% COMPONENTE 3: TEXTO CORRIDO E DESTAQUES
% -----------------------------------------------------------------------------
\newcommand{\termochave}[1]{\textcolor{BrandNavy}{\fontExtraBold #1}}
\newcommand{\destaqueazul}[1]{\textcolor{BrandSecondary}{\fontSemiBold #1}}
\newcommand{\codigoinline}[1]{%
  {\codeMonoFont\fontsize{9pt}{11pt}\selectfont\color{BrandSecondary}#1}%
}

\newcommand{\cenariopratico}[2]{%
  \par\vspace{6pt}\noindent
  {\fontsize{10pt}{13.5pt}\fontExtraBold\color{BrandNavy}#1}\par\nopagebreak
  \noindent #2\par\vspace{4pt}%
}
\newcommand{\blocoanalogia}[2]{\cenariopratico{#1}{#2}}

% -----------------------------------------------------------------------------
% COMPONENTE 4: FIGURAS COM LEGENDA NO TOPO E FONTE NA BASE
% -----------------------------------------------------------------------------
\newcounter{editorialfig}
\setcounter{editorialfig}{0}

% Uso: \figuracomfonte[<altura maxima>]{<arquivo>}{<largura maxima>}{<titulo>}{<fonte>}{<label>}
% O parametro opcional [<altura maxima>] controla o espaco vertical reservado
% a imagem (padrao 0.56\textheight). Diagramas verticais de grande porte
% (C4, Classes, Sequencia) exigem um valor maior para permanecerem legiveis.
\newcommand{\figuracomfonte}[6][0.56\textheight]{%
  \par\vspace{8pt}%
  \noindent\begin{minipage}{\textwidth}
    \refstepcounter{editorialfig}%
    \label{#6}%
    \addcontentsline{lof}{figure}{\protect\numberline{\theeditorialfig}#4}%
    \centering
    {\fontsize{9.5pt}{12.5pt}\selectfont\textcolor{BrandNavy}{\textbf{Figura \theeditorialfig}} \textendash\ \textcolor{BrandBodyText}{\fontSemiBold #4}}\par
    \vspace{5pt}%
    \includegraphics[max width=#3, max height=#1, keepaspectratio]{#2}\par
    \vspace{4pt}%
    {\fontsize{8.5pt}{11pt}\selectfont\raggedright\fontRegular\color{BrandGrayMed}\textbf{Fonte:} #5\par}%
  \end{minipage}%
  \par\vspace{8pt}%
}

\newcommand{\figuratikz}[4]{%
  \par\vspace{10pt}%
  \noindent\begin{minipage}{\textwidth}
    \refstepcounter{editorialfig}%
    \label{#3}%
    \addcontentsline{lof}{figure}{\protect\numberline{\theeditorialfig}#1}%
    \centering
    {\fontsize{9.5pt}{12.5pt}\selectfont\textcolor{BrandBodyText}{\fontExtraBold Figura \theeditorialfig} -- #1}\par
    \vspace{6pt}%
    #4\par
    \vspace{5pt}%
    {\fontsize{8.5pt}{11pt}\selectfont\raggedright\fontRegular\color{BrandGrayMed}\textbf{Fonte:} #2}\par
  \end{minipage}%
  \par\vspace{10pt}%
}

% -----------------------------------------------------------------------------
% COMPONENTE 5: QUADRO EDITORIAL E TABELA INSTITUCIONAL
% -----------------------------------------------------------------------------
\newcounter{editorialquadro}
\setcounter{editorialquadro}{0}

\newcommand{\quadroeditorial}[4]{%
  \par\vspace{10pt}%
  \noindent\begin{minipage}{\textwidth}
    \refstepcounter{editorialquadro}%
    \label{#3}%
    \addcontentsline{lot}{table}{\protect\numberline{Q\theeditorialquadro}#1}%
    \centering
    {\fontsize{9.5pt}{12.5pt}\selectfont\textcolor{BrandBodyText}{\fontExtraBold Quadro \theeditorialquadro} -- #1}\par
    \vspace{4pt}%
    \begin{tcolorbox}[
      enhanced,
      colback=white,
      colframe=BrandGrayRule,
      boxrule=0.9pt,
      arc=4pt,
      left=12pt,
      right=12pt,
      top=10pt,
      bottom=10pt,
      fontupper=\fontsize{9.2pt}{13pt}\selectfont\color{BrandBodyText}
    ]
      #4
    \end{tcolorbox}
    \vspace{2pt}%
    {\fontsize{8.5pt}{11pt}\fontRegular\color{BrandGrayMed}Fonte: #2}\par
  \end{minipage}%
  \par\vspace{10pt}%
}

\newcounter{editorialtabela}
\setcounter{editorialtabela}{0}

\newcommand{\fonte}[1]{%
  \par\vspace{3pt}%
  {\fontsize{8.5pt}{11pt}\selectfont\raggedright\fontRegular\color{BrandGrayMed}\textbf{Fonte:} #1\par}%
}

\newcommand{\tabelaabnt}[4]{%
  \par\vspace{6pt}%
  \noindent\begin{minipage}{\textwidth}
    \refstepcounter{editorialtabela}%
    \label{#3}%
    \addcontentsline{lot}{table}{\protect\numberline{\theeditorialtabela}#1}%
    \centering
    {\fontsize{9.5pt}{12pt}\selectfont\textcolor{BrandNavy}{\textbf{Tabela \theeditorialtabela}} \textendash\ \textcolor{BrandBodyText}{\fontSemiBold #1}}\par
    \vspace{4pt}%
    \begingroup
      \renewcommand{\tabularxcolumn}[1]{m{##1}}%
      \renewcommand{\arraystretch}{1.20}%
      \setlength{\tabcolsep}{4.5pt}%
      \emergencystretch=2em%
      \tolerance=2500%
      \hyphenpenalty=50%
      \hbadness=2500%
      #4\par
    \endgroup
    \vspace{3pt}%
    {\fontsize{8.5pt}{11pt}\selectfont\raggedright\fontRegular\color{BrandGrayMed}\textbf{Fonte:} #2\par}%
  \end{minipage}%
  \par\vspace{8pt}%
}

\newcommand{\tabelaeditorial}[5]{%
  \par\vspace{10pt}%
  \noindent\begin{minipage}{\textwidth}
    \refstepcounter{editorialtabela}%
    \label{#3}%
    \addcontentsline{lot}{table}{\protect\numberline{\theeditorialtabela}#1}%
    \centering
    {\fontsize{9.5pt}{12.5pt}\selectfont\textcolor{BrandBodyText}{\fontExtraBold Tabela \theeditorialtabela} -- #1}\par
    \vspace{6pt}%
    {\fontsize{9pt}{12.5pt}\selectfont
     \renewcommand{\arraystretch}{1.30}%
     \setlength{\tabcolsep}{7pt}%
     \begin{tabularx}{\textwidth}{#4}
       \arrayrulecolor{BrandNavy}\toprule[1.2pt]
       #5
       \arrayrulecolor{BrandNavy}\bottomrule[1.2pt]
     \end{tabularx}%
    }\par
    \vspace{5pt}%
    {\fontsize{8.5pt}{11pt}\fontRegular\color{BrandGrayMed}Fonte: #2}\par
  \end{minipage}%
  \par\vspace{10pt}%
}

% -----------------------------------------------------------------------------
% COMPONENTE 6: VERIFICAÇÃO DE CONFORMIDADE
% -----------------------------------------------------------------------------
\newcounter{editorialcheck}
\setcounter{editorialcheck}{0}

\newcommand{\verificacaoconformidade}[2]{%
  \refstepcounter{editorialcheck}%
  \par\vspace{10pt}%
  \begin{tcolorbox}[
    enhanced,
    breakable,
    colback=BrandTableAlt,
    colframe=BrandNavy,
    boxrule=0.8pt,
    leftrule=3.5pt,
    arc=2pt,
    left=12pt,
    right=12pt,
    top=10pt,
    bottom=10pt,
    title={\fontsize{9.5pt}{12pt}\fontExtraBold\color{white}%
           VERIFICAÇÃO DE CONFORMIDADE \theeditorialcheck\ — #1},
    colbacktitle=BrandNavy,
    toptitle=5pt,
    bottomtitle=5pt,
    fontupper=\fontsize{9.2pt}{13pt}\selectfont\color{BrandBodyText}
  ]
    #2
  \end{tcolorbox}%
  \par\vspace{8pt}%
}

\newcommand{\exerciciofixacao}[2]{\verificacaoconformidade{#1}{#2}}

% -----------------------------------------------------------------------------
% COMPONENTE 7: CAIXA CURIOSIDADES / NOTA TÉCNICA (DEEP SAPPHIRE SÓBRIO)
% -----------------------------------------------------------------------------
\newcommand{\caixacuriosidades}[2][NOTA TÉCNICA DE ENGENHARIA]{%
  \par\vspace{12pt}%
  \begin{tcolorbox}[
    enhanced,
    breakable,
    colback=BoxSapphireBg,
    colframe=BoxSapphireBg,
    boxrule=0pt,
    arc=2pt,
    left=14pt,
    right=14pt,
    top=12pt,
    bottom=12pt,
    fontupper=\fontsize{9.5pt}{13.5pt}\selectfont\color{white}
  ]
    {\fontsize{11pt}{14pt}\fontExtraBold\color{white}\MakeUppercase{#1}}\par
    \vspace{3pt}%
    \textcolor{white!30}{\rule{3.0cm}{0.6pt}}\par
    \vspace{5pt}%
    #2
  \end{tcolorbox}%
  \par\vspace{10pt}%
}

\newcommand{\notatecnica}[2][NOTA TÉCNICA DE ENGENHARIA]{%
  \caixacuriosidades[#1]{#2}%
}

% -----------------------------------------------------------------------------
% COMPONENTE 8: CAIXA SAIBA MAIS (CARMIM SÓBRIO #7A1D2E)
% -----------------------------------------------------------------------------
\newcommand{\caixasaibamais}[2][SAIBA MAIS — NORMAS E APROFUNDAMENTO]{%
  \par\vspace{12pt}%
  \begin{tcolorbox}[
    enhanced,
    breakable,
    colback=BoxCrimsonBg,
    colframe=BoxCrimsonBg,
    boxrule=0pt,
    arc=2pt,
    left=14pt,
    right=14pt,
    top=12pt,
    bottom=12pt,
    fontupper=\fontsize{9.5pt}{13.5pt}\selectfont\color{white}
  ]
    {\fontsize{11pt}{14pt}\fontExtraBold\color{white}\MakeUppercase{#1}}\par
    \vspace{3pt}%
    \textcolor{white!30}{\rule{3.0cm}{0.6pt}}\par
    \vspace{5pt}%
    #2
  \end{tcolorbox}%
  \par\vspace{10pt}%
}

\newcommand{\saibamaisnormativo}[2][SAIBA MAIS — NORMAS E APROFUNDAMENTO]{%
  \caixasaibamais[#1]{#2}%
}

% -----------------------------------------------------------------------------
% COMPONENTE 9: CAIXA PARA RELEMBRAR (CARVÃO MEIA-NOITE #161C2B)
% -----------------------------------------------------------------------------
\newcommand{\caixarelembrar}[2][SÍNTESE EXECUTIVA DA UNIDADE]{%
  \par\vspace{12pt}%
  \begin{tcolorbox}[
    enhanced,
    breakable,
    colback=BoxMidnightBg,
    colframe=BoxMidnightBg,
    boxrule=0pt,
    arc=2pt,
    left=14pt,
    right=14pt,
    top=12pt,
    bottom=12pt,
    fontupper=\fontsize{9.5pt}{13.5pt}\selectfont\color{white}
  ]
    {\fontsize{11pt}{14pt}\fontExtraBold\color{white}\MakeUppercase{#1}}\par
    \vspace{3pt}%
    \textcolor{white!30}{\rule{3.0cm}{0.6pt}}\par
    \vspace{5pt}%
    #2
  \end{tcolorbox}%
  \par\vspace{10pt}%
}

\newcommand{\sintesecapitulo}[2][SÍNTESE EXECUTIVA DA UNIDADE]{%
  \caixarelembrar[#1]{#2}%
}

% -----------------------------------------------------------------------------
% COMPONENTE 10: LISTAS, AVISO E CÓDIGO
% -----------------------------------------------------------------------------
\newenvironment{listapontos}{%
  \begin{itemize}[
    label={\textcolor{BrandCrimson}{\fontsize{10pt}{10pt}\fontExtraBold »}},
    leftmargin=0.95cm,
    labelsep=0.25cm,
    itemsep=3pt,
    topsep=4pt,
    parsep=2pt
  ]
}{%
  \end{itemize}%
}

\newcommand{\itemponto}[2]{%
  \item \textcolor{BrandNavy}{\fontExtraBold #1:} #2%
}

\newenvironment{listaetapas}{%
  \begin{enumerate}[
    label={\textcolor{BrandCrimson}{\fontExtraBold\arabic*.}},
    leftmargin=0.95cm,
    labelsep=0.25cm,
    itemsep=4pt,
    topsep=5pt,
    parsep=2pt
  ]
}{%
  \end{enumerate}%
}

\newcommand{\itemetapa}[2]{%
  \item \textcolor{BrandNavy}{\fontExtraBold #1:} #2%
}

\newcommand{\caixaatencao}[2][Atenção:]{%
  \par\vspace{8pt}%
  \begin{tcolorbox}[
    enhanced,
    colback=BrandTableAlt,
    colframe=BrandSecondary,
    boxrule=0.8pt,
    leftrule=3.0pt,
    arc=2pt,
    left=12pt,
    right=12pt,
    top=8pt,
    bottom=8pt,
    fontupper=\fontsize{9.2pt}{13pt}\selectfont\color{BrandBodyText}
  ]
    {\fontsize{9.8pt}{13pt}\fontExtraBold\color{BrandSecondary}#1}\par\vspace{2pt}%
    #2
  \end{tcolorbox}%
  \par\vspace{8pt}%
}

% Bloco de Código Estilo Dark Sóbrio
\newcounter{editorialcodigo}
\setcounter{editorialcodigo}{0}

\newtcblisting{editoriallisting}[4][]{%
  enhanced,
  colback=BrandCodeBg,
  colframe=BrandCodeBg,
  arc=3pt,
  boxrule=0pt,
  left=10pt,
  right=10pt,
  top=8pt,
  bottom=10pt,
  overlay={
    \node[
      fill=BrandCodePill,
      text=white,
      font=\fontsize{7.5pt}{9pt}\fontSemiBold,
      inner xsep=6pt,
      inner ysep=2.5pt,
      rounded corners=2pt,
      anchor=south east
    ] at ($(frame.south east) + (-8pt, 5pt)$) {#3};
  },
  listing only,
  listing options={
    basicstyle=\codeMonoFont\fontsize{8.4pt}{11.2pt}\selectfont\color{BrandCodeText},
    commentstyle=\color{BrandCodeComment}\itshape,
    keywordstyle=\color{BrandCodeKeyword}\bfseries,
    stringstyle=\color{BrandCodeString},
    breaklines=true,
    breakatwhitespace=true,
    tabsize=2,
    showstringspaces=false,
    columns=fullflexible,
    keepspaces=true,
    frame=none
  },
  #1
}

\newcommand{\blocodigocaption}[3]{%
  \refstepcounter{editorialcodigo}%
  \label{#3}%
  \centering
  {\fontsize{9.5pt}{12.5pt}\selectfont\textcolor{BrandBodyText}{\fontExtraBold Código \theeditorialcodigo} -- #1}\par
  \vspace{4pt}%
}

\newcommand{\blocodigofonte}[1]{%
  \vspace{2pt}%
  \centering
  {\fontsize{8.5pt}{11pt}\fontRegular\color{BrandGrayMed}Fonte: #1}\par
  \vspace{8pt}%
}

% -----------------------------------------------------------------------------
% COMPONENTE 11: NOTA DE RODAPÉ
% -----------------------------------------------------------------------------
\newcommand{\notarodape}[1]{\footnote{#1}}

% -----------------------------------------------------------------------------
% COMPONENTE 12: REFERÊNCIAS
% -----------------------------------------------------------------------------
\newcommand{\refitem}[1]{%
  \par\noindent
  \hangindent=0cm\hangafter=1
  {\fontsize{9.5pt}{13.5pt}\selectfont\color{BrandBodyText}#1}\par\vspace{7pt}%
}

\newcommand{\subsecaoreferencias}[1]{%
  \par\vspace{10pt}\noindent
  {\fontsize{11.5pt}{15pt}\fontExtraBold\color{BrandNavy}#1}\par\vspace{5pt}%
}

\newcommand{\refsite}[3]{%
  \par\noindent
  {\fontsize{9.5pt}{13.5pt}\selectfont\color{BrandBodyText}%
   #1\ Disponível em: \href{#2}{\textcolor{BrandCyanLink}{\underline{#2}}}. Acesso em: #3.}\par\vspace{7pt}%
}
